\chapter{Fontes de dados}\label{anx:data}

Todos os conjuntos de dados provêm do catálogo público da plataforma de computação. Eles são acessados
em resolução nativa e reescalonados para unidades físicas quando o provedor armazena as bandas como
inteiros; cada fonte foi conferida com o catálogo ativo no momento do processamento.

\begin{table}[htbp]
\centering
\caption{Fontes de dados biofísicos e agroclimáticos.}
\label{tab:src-biophys}
\begin{tabular}{@{}>{\raggedright\arraybackslash}p{2.7cm}>{\raggedright\arraybackslash}p{3.6cm}>{\raggedright\arraybackslash}p{2cm}>{\raggedright\arraybackslash}p{5.5cm}@{}}
\toprule
Tema & Conjunto de dados (provedor) & Resolução nativa & Variáveis utilizadas \\
\midrule
Clima e balanço hídrico & TerraClimate, mensal & $\approx$4,6~km & Precipitação,
  evapotranspiração potencial e real, déficit hídrico climático, umidade do solo, déficit de
  pressão de vapor, radiação solar, temperatura máxima e mínima \\
Terreno & Modelo digital de elevação SRTM & 30~m & Elevação, declividade, orientação, posição
  topográfica \\
Hidrografia & Acumulação de fluxo HydroSHEDS & 15 segundos de arco & Índice de umidade
  topográfica, densidade de drenagem \\
Solo (0--30~cm) & Propriedades de solo OpenLandMap & 250~m & Argila, areia, carbono orgânico,
  pH (H$_2$O), densidade aparente, teor de água disponível para as plantas \\
Água de superfície & JRC Global Surface Water & 30~m & Ocorrência e sazonalidade
  $\rightarrow$ distância à água permanente \\
Cobertura da terra & MapBiomas, Coleção~10 (2024) & 30~m & Máscara e frações por classe
  (apenas máscara e contexto) \\
Áreas protegidas & World Database on Protected Areas (WDPA) & vetor & Distância a reservas
  $\rightarrow$ conectividade de conservação (camada de governança, e não de cobertura da
  terra) \\
Biomassa aérea & Densidade global de carbono na biomassa aérea, ano de
  referência 2010 \citep[p.~1]{spawn_harmonized_2020} & $\approx$300~m & Valor do serviço de
  armazenamento de carbono (conservação) \\
Clima futuro & NASA NEX-GDDP-CMIP6, diário, com redução estatística de escala (\emph{downscaling}) &
  $\approx$25~km & Precipitação, temperatura máxima e mínima $\rightarrow$ fatores de mudança
  (\emph{delta-change}) \\
\bottomrule
\end{tabular}
\end{table}

As bandas de solo são a média dos horizontes de 0, 10 e 30~cm. As conversões de unidade verificadas
são: carbono orgânico do solo $\times 5$ (média territorial de $\approx$13~g~kg$^{-1}$), pH
$\times 0{,}1$ ($\approx$4,7 a 7,3) e densidade aparente $\times 10$ ($\approx$1.259~kg~m$^{-3}$);
argila, areia e teor de água não requerem reescalonamento.

\begin{table}[htbp]
\centering
\caption{Fontes de dados de uso atual e validação.}
\label{tab:src-realized}
\begin{tabular}{@{}>{\raggedright\arraybackslash}p{5cm}>{\raggedright\arraybackslash}p{2.7cm}>{\raggedright\arraybackslash}p{6cm}@{}}
\toprule
Fonte (provedor) & Resolução nativa & Finalidade \\
\midrule
Malha Municipal Digital do IBGE (2025) & $\approx$111~m (geometria simplificada) & Agregação e
  apresentação de resultados por município (247 municípios; sem \emph{upload}) \\
Uso e cobertura da terra anual MapBiomas Brasil & 30~m & Máscaras, análise de potencial
  \emph{versus} uso atual, caracterização descritiva da composição \\
Produtividade primária líquida (anual) e bruta (8 dias) do MODIS & 500~m & Validação da
  aptidão (substituto de rendimento) \\
\bottomrule
\end{tabular}
\end{table}

\chapter{Especificação completa de aptidão e pesos AHP}\label{anx:spec}

A Tabela~\ref{tab:spec-longtable} reproduz a especificação completa por segmento que fundamenta a
Seção~\ref{sec:suitability}. A declividade está em graus. Os limiares são valores de primeira
aproximação ancorados na literatura e ajustados, pela recalibração empírica, à faixa observada de cada
variável em GO~e~DF. Os pesos, por sua vez, são ajustados ao poder discriminante amostral de cada
fator (Seção~\ref{sec:ahp}). Assim, um fator quase constante em GO~e~DF recebe peso baixo, por mais
importante que a literatura o considere.

\begin{longtable}{@{}>{\raggedright\arraybackslash}p{2.6cm}>{\raggedright\arraybackslash}p{3.6cm}p{0.9cm}>{\raggedright\arraybackslash}p{1.9cm}>{\raggedright\arraybackslash}p{3.2cm}@{}}
\caption{Fatores, pesos, formas de pertinência e limiares por segmento.}\label{tab:spec-longtable} \\
\toprule
Segmento & Fator & Peso & Forma & Limiares \\
\midrule
\endfirsthead
\toprule
Segmento & Fator & Peso & Forma & Limiares \\
\midrule
\endhead
% GENERATED by tools/gen_thesis_tables.py -- do not edit by hand.
% Edit the source CSV and regenerate; see CLAUDE.md §9.
\textbf{Soja} & Declive ($^\circ$) & 0,42 & decrescente & 4, 8 \\
 & pH do solo & 0,33 & intervalar & 4,5, 5,5, 6,8, 7,5 \\
 & Teor de argila (\%) & 0,19 & intervalar & 20, 32, 45, 62 \\
 & Temp. do trimestre mais quente ($^\circ$C) & 0,06 & intervalar & 18, 22, 30, 34 \\
 & Precipitação anual (mm) & 0,00 & intervalar & 800, 1100, 1800, 2200 \\
 & Índice de aridez (P/ETP) & 0,00 & intervalar & 0,4, 0,7, 1,3, 1,6 \\
\textbf{Cana-de-açúcar} & Teor de água disponível para as plantas & 0,30 & crescente & 22,1, 33,6 \\
 & Declive ($^\circ$) & 0,29 & decrescente & 5, 10 \\
 & Teor de argila (\%) & 0,21 & intervalar & 22, 32, 52, 66 \\
 & Precipitação anual (mm) & 0,15 & intervalar & 900, 1300, 2000, 2600 \\
 & Meses secos (n) & 0,05 & intervalar & 3, 5, 5, 9 \\
 & Temp. do trimestre mais quente ($^\circ$C) & 0,00 & intervalar & 20, 23, 32, 36 \\
\textbf{Outras culturas anuais} & Carbono orgânico do solo (g/kg) & 0,48 & crescente & 6, 20 \\
 & Graus-dia de crescimento & 0,23 & intervalar & 2500, 3500, 5500, 7000 \\
 & pH do solo & 0,14 & intervalar & 4,5, 5,5, 6,8, 7,5 \\
 & Declive ($^\circ$) & 0,11 & decrescente & 4, 8 \\
 & CV da precipitação & 0,04 & decrescente & 0,6, 1,3 \\
 & Precipitação anual (mm) & 0,00 & intervalar & 700, 1000, 1700, 2200 \\
\textbf{Piscicultura} & Distância à água sazonal (m) & 0,49 & decrescente & 750, 6000 \\
 & Temp. do trimestre mais quente ($^\circ$C) & 0,20 & intervalar & 22, 26, 32, 36 \\
 & Teor de argila (\%) & 0,14 & crescente & 15, 35 \\
 & Declive ($^\circ$) & 0,11 & decrescente & 2,86, 5,71 \\
 & Densidade de drenagem & 0,05 & crescente & 0, 0,09 \\
\textbf{Pecuária} & Umidade do solo (mm) & 0,65 & crescente & 51, 120 \\
 & Declive ($^\circ$) & 0,15 & decrescente & 6, 15 \\
 & Carbono orgânico do solo (g/kg) & 0,10 & crescente & 8, 17 \\
 & Meses secos (n) & 0,09 & decrescente & 5, 9 \\
 & Precipitação anual (mm) & 0,00 & intervalar & 700, 1000, 2000, 2600 \\
\textbf{Conservação} & Distância a área protegida (m) & 0,34 & decrescente & 5000, 40000 \\
 & Carbono em biomassa acima do solo (Mg~C/ha) & 0,17 & crescente & 10, 20 \\
 & Rugosidade do terreno (m) & 0,17 & crescente & 5, 30 \\
 & Índice de umidade topográfica & 0,07 & crescente & 9,44, 11,56 \\
 & Índice de aridez (P/ETP) $\dagger$ & 0,07 & crescente & 0,95, 1,25 \\
 & Temp. do trimestre mais quente ($^\circ$C) $\dagger$ & 0,06 & decrescente & 23, 31 \\
 & Fração de vegetação nativa & 0,06 & crescente & 0,25, 0,85 \\
 & Densidade de drenagem & 0,03 & crescente & 0,01, 0,08 \\
 & Declive ($^\circ$) & 0,02 & crescente & 2,86, 11,31 \\
\textbf{Geração fotovoltaica} & Radiação solar (W/m$^2$) & 0,44 & crescente & 186, 204 \\
 & Temp. do trimestre mais quente ($^\circ$C) & 0,25 & decrescente & 24, 31 \\
 & Índice de claridade (cobertura de nuvens) & 0,16 & crescente & 0,47, 0,5 \\
 & Declive ($^\circ$) & 0,09 & decrescente & 2,86, 8,53 \\
 & Componente norte da orientação & 0,06 & crescente & -0,3, 0,5 \\
\bottomrule
   % the fragment ends with \bottomrule
\end{longtable}

A conservação ($\dagger$) inclui dois fatores climáticos \textbf{sensíveis} (aridez e
temperatura do trimestre mais quente), que a tornam responsiva ao CMIP6, e é agregada pela
média aritmética ponderada \textbf{compensatória} (Seção~\ref{sec:aggregation}). Sua fração de
vegetação nativa (peso 0,06) é a única entrada de cobertura da terra. Ela recebe peso reduzido e é
admitida pela salvaguarda contra circularidade. Os outros seis segmentos usam a média geométrica de
fator limitante.

Nenhum fator é retirado da matriz AHP. Sete fatores climáticos têm $\sigma(\mu) = 0$ no presente e
permanecem na matriz graças à ponderação sobre o horizonte CMIP6 descrita na Seção~\ref{sec:ahp};
cinco deles ainda recebem peso exatamente nulo, por terem $\mu \equiv 1$ nas duas populações. A coluna
\emph{Fatores} da Tabela~\ref{tab:ahp-cr} conta, portanto, todos os fatores de cada segmento,
inclusive os de peso nulo.

Os pesos de cada segmento são as prioridades de uma matriz de Saaty de comparação par a par
documentada, e todas as matrizes são amplamente consistentes, nos termos definidos na
Seção~\ref{sec:ahp} (Tabela~\ref{tab:ahp-cr}).

\begin{table}[htbp]
\centering
\caption{Razão de consistência AHP por segmento.}
\label{tab:ahp-cr}
% GENERATED by tools/gen_thesis_tables.py -- do not edit by hand.
% Edit the source CSV and regenerate; see CLAUDE.md §9.
\begin{tabular}{@{}lrr@{}}
\toprule
Segmento & Fatores ($n$) & Razão de consistência \\
\midrule
Soja & 6 & 0,0133 \\
Cana-de-açúcar & 6 & 0,0335 \\
Outras culturas anuais & 6 & 0,0054 \\
Piscicultura & 5 & 0,0168 \\
Pecuária & 5 & 0,0067 \\
Conservação & 9 & 0,0245 \\
Geração fotovoltaica & 5 & 0,0205 \\
\bottomrule
\end{tabular}

\end{table}

\chapter{Derivação dos atributos}\label{anx:derivations}

\textbf{Clima e balanço hídrico (14 variáveis).} Calcula-se uma climatologia de doze meses
para cada banda climática sobre o registro de 1991--2020 (a média climatológica de cada mês
do calendário); os agregados anuais são a soma das doze médias mensais. As variáveis
derivadas do balanço hídrico são a precipitação anual, a evapotranspiração potencial anual, o
déficit hídrico climático, a evapotranspiração real e o índice de aridez (precipitação anual
$\div$ evapotranspiração potencial anual). As variáveis de sazonalidade são:
\begin{itemize}
  \item a sazonalidade da precipitação (o coeficiente de variação das doze médias mensais);
  \item a duração da estação seca (meses com precipitação média abaixo de 60~mm);
  \item a precipitação da estação chuvosa (a soma sobre os meses mais úmidos que a média mensal geral).
\end{itemize}
Completam o conjunto: a umidade média do solo; o déficit de pressão de vapor e a radiação solar
médios; os graus-dia de crescimento; e a temperatura média do trimestre mais quente e do mais frio (o
máximo e o mínimo das doze médias móveis de três meses). Os graus-dia de crescimento são acumulados a
partir da climatologia mensal de temperatura média, com temperatura base
$T_{\text{base}} = 10\,^\circ\mathrm{C}$:
\begin{equation*}
  \mathrm{GDD} = \sum_{m=1}^{12} \max\!\left(\bar{T}_m - T_{\text{base}},\, 0\right)\, N_m,
\end{equation*}
em que $\bar{T}_m$ é a temperatura média do mês $m$ e $N_m = 30{,}4$ dias é a duração média do mês. É o
Método~1 de \citet[p.~292]{mcmaster_growing_1997} aplicado a médias mensais, o que tende a
subestimar o GDD em meses próximos da temperatura base.

\textbf{Terreno (6 variáveis).} A elevação e a declividade vêm do MDE de 30~m. A orientação é
decomposta em uma componente norte ($\cos$) e uma leste ($\sin$), para que a variável circular
possa ser usada em aritmética de pixels. O índice de posição topográfica é a diferença entre
a elevação $z_0$ de uma célula e a elevação média de sua vizinhança circular,
$\mathrm{TPI} = z_0 - \overline{z}_{\text{focal}}$ (positivo em cristas, negativo em vales). O
índice de umidade topográfica é $\mathrm{TWI} = \ln(\alpha / \tan\beta)$
\citep[p.~48, eq.~8]{beven_physically_1979}, em que $\beta$ é a declividade local e $\alpha$, a área
de contribuição específica. Essa área é aproximada por $(n+1)\,\Delta x$, em que $n$ é a acumulação
de fluxo do HydroSHEDS em número de células e $\Delta x \approx 464$~m é a largura da sua célula de
15 segundos de arco.

\textbf{Solo (6 variáveis).} Usa-se a média de 0--30~cm (horizontes de 0, 10 e 30~cm) para argila,
areia, carbono orgânico, pH (H$_2$O), densidade aparente e teor de água disponível para as
plantas, cada um reescalonado para unidades físicas pelo seu fator de conversão verificado
(Anexo~\ref{anx:data}).

\textbf{Água (3 variáveis).} Uma máscara de água permanente (ocorrência de água superficial
$>$~80\%) produz uma superfície de distância euclidiana à água na grade de 250~m. A sazonalidade da
água de superfície é obtida diretamente do produto de origem. A densidade de drenagem é a fração de
células de canal (acumulação de fluxo acima de um limiar fixo) em uma vizinhança circular; sua faixa
empírica, de $\approx$0--0,15 (percentil 90 $\approx$ 0,095), define seus limiares de pertinência.

\textbf{Local (2 variáveis; fora da concatenação, ver Seção~\ref{sec:zoning}).} O \emph{índice de
claridade} atmosférica é a razão entre a radiação solar observada na superfície e a radiação
extraterrestre no topo da atmosfera, $K_t = R_s/R_a$, truncada a $[0;1]$. $R_a$ é calculada
analiticamente pela formulação FAO-56, apenas como função da latitude e do dia do ano (declinação
solar, ângulo horário do pôr do sol e inverso da distância relativa Terra--Sol), integrada mês a mês
e tomada em média anual. A normalização é o que importa: $R_s$ isolada é quase uniforme em GO~e~DF,
enquanto $K_t$ mede a fração da radiação disponível que de fato chega ao solo e, por isso, carrega o
sinal espacial de nebulosidade. A segunda variável é a distância euclidiana à água superficial
\textbf{semipermanente} (ocorrência $>$~50\% no JRC Global Surface Water, isto é, água presente cerca
de seis meses ou mais por ano), calculada pela mesma transformada de distância usada para a água
permanente. Ela existe porque o critério de local para aquicultura em pequenos reservatórios depende da
água sazonal, e não da permanente. Quando ancorado na água permanente, o fator de piscicultura
apresentava distâncias médias de cerca de 13~km e permanecia no valor mínimo de pertinência.

\textbf{Conservação (3 variáveis; fora da concatenação, ver Seção~\ref{sec:zoning}).} A
\emph{distância a áreas protegidas} rasteriza os polígonos do WDPA e aplica a mesma transformada de
distância (zero dentro de uma reserva, crescendo para fora). Ela entra com pertinência decrescente,
de modo que a vizinhança de reservas recebe pontuação alta. Trata-se de um \emph{proxy} de
conectividade e consolidação; o WDPA é uma camada de governança, e não de cobertura da terra. A
\emph{rugosidade do terreno} é o desvio-padrão da elevação em uma vizinhança circular de 3 pixels
($\approx$750~m) sobre a grade de 250~m, lida como heterogeneidade de micro-hábitat e, ao mesmo tempo,
como baixa aptidão ao cultivo mecanizado. O \emph{carbono da biomassa aérea} é a densidade de carbono
na biomassa lenhosa viva \citep{spawn_harmonized_2020}, reamostrada bilinearmente para a grade, e
representa o valor do serviço ecossistêmico de estoque de carbono. As duas últimas têm pertinência
crescente. A correlação moderada do carbono com a cobertura florestal é uma ressalva declarada, e é
por isso que ele recebe peso modesto (Tabela~\ref{tab:spec-longtable}).

\textbf{Máscara de cobertura da terra (apenas contexto).} O MapBiomas, a 30~m, é
agregado a 250~m como uma classe majoritária e frações por classe. Disso resultam uma máscara de
áreas \emph{excluídas} (urbano, água permanente, área úmida herbácea), uma máscara de áreas
\emph{disponíveis} (lavoura, campo, vegetação arbustiva, solo exposto) e frações de
árvore, lavoura e campo para contexto.

\chapter{Tabelas de resultados suplementares}\label{anx:tables}

Este anexo reúne as tabelas de resultados citadas no Capítulo~\ref{ch:results}. Na
Tabela~\ref{tab:spearman-npp}, o \emph{proxy} de NPP anual é lido apenas como \emph{contexto} de
produtividade, pelas razões dadas na Seção~\ref{sec:validation-present}; os validadores primários das
culturas são as métricas de discriminação por presença (Tabela~\ref{tab:presence-auc}).

\begin{table}[htbp]
\centering
\caption{Médias biofísicas discriminantes por zona.}
\label{tab:zone-means}
% GENERATED by tools/gen_thesis_tables.py -- do not edit by hand.
% Edit the source CSV and regenerate; see CLAUDE.md §9.
\begin{tabular}{@{}lrrrr@{}}
\toprule
Variável & Z1 & Z2 & Z3 & Z4 \\
\midrule
Tamanho da amostra (pontos) & 10.163 & 2.768 & 6.089 & 7.800 \\
Parcela da área de interesse & 46,4\% & 10,4\% & 15,2\% & 28,1\% \\
Elevação (m) & 784 & 733 & 483 & 445 \\
Declividade ($^\circ$) & 4,5 & 16,9 & 3,5 & 3,3 \\
Índice de umidade topográfica & 10,2 & 8,8 & 11,4 & 10,6 \\
Temp. do trimestre mais quente ($^\circ$C) & 24,5 & 25,5 & 25,6 & 26,6 \\
Índice de aridez (P/ETP) & 1,15 & 1,10 & 1,13 & 1,05 \\
Umidade do solo (mm) & 86 & 110 & 93 & 109 \\
Teor de argila (\%) & 37,0 & 34,4 & 29,4 & 26,4 \\
Carbono orgânico do solo (g/kg) & 15,0 & 16,0 & 13,5 & 11,1 \\
Densidade de drenagem (fração) & 0,006 & 0,009 & 0,107 & 0,012 \\
Distância à água (m) & 16.181 & 19.624 & 8.510 & 15.620 \\
\bottomrule
\end{tabular}

\end{table}

\begin{table}[htbp]
\centering
\caption{Mudança média projetada, $\Delta S$, por zona para as três culturas sensíveis ao clima,
nas quatro combinações de cenário e janela.}
\label{tab:zone-shift}
% GENERATED by tools/gen_thesis_tables.py -- do not edit by hand.
% Edit the source CSV and regenerate; see CLAUDE.md §9.
\begin{tabular}{@{}llrrrr@{}}
\toprule
 &  & \multicolumn{2}{c}{SSP2-4.5} & \multicolumn{2}{c}{SSP5-8.5} \\
\cmidrule(lr){3-4} \cmidrule(lr){5-6}
Segmento & Zona & 2031--50 & 2051--70 & 2031--50 & 2051--70 \\
\midrule
Soja & Z1 & 0,000 & 0,000 & 0,000 & 0,000 \\
Soja & Z2 & 0,000 & 0,000 & 0,000 & 0,000 \\
Soja & Z3 & 0,000 & 0,000 & 0,000 & $-$0,001 \\
Soja & Z4 & 0,000 & 0,000 & 0,000 & $-$0,004 \\
\midrule
Cana-de-açúcar & Z1 & 0,000 & 0,000 & $-$0,001 & 0,000 \\
Cana-de-açúcar & Z2 & 0,000 & 0,000 & 0,000 & 0,000 \\
Cana-de-açúcar & Z3 & 0,000 & 0,000 & 0,000 & 0,000 \\
Cana-de-açúcar & Z4 & $-$0,001 & $-$0,001 & $-$0,001 & $-$0,002 \\
\midrule
Outras culturas anuais & Z1 & $-$0,004 & $-$0,012 & $-$0,008 & $-$0,036 \\
Outras culturas anuais & Z2 & $-$0,003 & $-$0,006 & $-$0,005 & $-$0,012 \\
Outras culturas anuais & Z3 & $-$0,012 & $-$0,025 & $-$0,018 & $-$0,053 \\
Outras culturas anuais & Z4 & $-$0,027 & $-$0,048 & $-$0,038 & $-$0,093 \\
\bottomrule
\end{tabular}

\end{table}

\begin{table}[htbp]
\centering
\caption{Parcela da terra avaliada de cada segmento rebaixada em pelo menos uma classe da
FAO, presente $\rightarrow$ futuro (as promoções, desprezíveis, foram omitidas).}
\label{tab:downgrade-share}
\small
% GENERATED by tools/gen_thesis_tables.py -- do not edit by hand.
% Edit the source CSV and regenerate; see CLAUDE.md §9.
\begin{tabular}{@{}lrrrr@{}}
\toprule
 & \multicolumn{2}{c}{SSP2-4.5} & \multicolumn{2}{c}{SSP5-8.5} \\
\cmidrule(lr){2-3} \cmidrule(lr){4-5}
Segmento & 2031--50 & 2051--70 & 2031--50 & 2051--70 \\
\midrule
Soja & 0,0\% & 0,1\% & 0,0\% & 0,8\% \\
Cana-de-açúcar & 0,1\% & 0,1\% & 0,3\% & 0,2\% \\
Outras culturas anuais & 4,8\% & 9,4\% & 7,0\% & 21,8\% \\
Piscicultura & 0,0\% & 0,0\% & 0,0\% & 0,0\% \\
Pecuária & 0,0\% & 0,0\% & 0,0\% & 0,5\% \\
Conservação & 14,4\% & 16,7\% & 17,2\% & 25,4\% \\
Geração fotovoltaica & 11,2\% & 18,0\% & 15,0\% & 32,9\% \\
\bottomrule
\end{tabular}

\end{table}

\begin{table}[htbp]
\centering
\caption{Matriz de transição de classes da FAO, presente $\rightarrow$ futuro, para as
outras culturas anuais sob SSP5-8.5, 2051--2070 (área como percentual da terra avaliada do
segmento).}
\label{tab:transition-othercrops}
% GENERATED by tools/gen_thesis_tables.py -- do not edit by hand.
% Edit the source CSV and regenerate; see CLAUDE.md §9.
\begin{tabular}{@{}lrrrrr@{}}
\toprule
Presente $\downarrow$ / futuro $\rightarrow$ & N & S3 & S2 & S1 & Total \\
\midrule
N & 36,0\% & 0,0\% & 0,0\% & 0,0\% & \textbf{36,0\%} \\
S3 & 3,7\% & 15,2\% & 0,0\% & 0,0\% & \textbf{18,9\%} \\
S2 & 0,2\% & 12,7\% & 19,8\% & 0,0\% & \textbf{32,6\%} \\
S1 & 0,0\% & 0,1\% & 5,1\% & 7,3\% & \textbf{12,4\%} \\
\bottomrule
\end{tabular}

\end{table}

\begin{table}[htbp]
\centering
\caption{Composição do uso atual da terra dentro de cada zona biofísica (parcela da área da
zona).}
\label{tab:zone-realized}
% GENERATED by tools/gen_thesis_tables.py -- do not edit by hand.
% Edit the source CSV and regenerate; see CLAUDE.md §9.
\begin{tabular}{@{}lrrrr@{}}
\toprule
Uso atual & Z1 & Z2 & Z3 & Z4 \\
\midrule
Soja & 23,1\% & 0,1\% & 7,9\% & 5,4\% \\
Cana-de-açúcar & 3,5\% & 0,0\% & 3,1\% & 0,6\% \\
Outras culturas anuais & 0,8\% & 0,0\% & 1,0\% & 1,3\% \\
Pastagem & 44,0\% & 29,7\% & 58,1\% & 66,5\% \\
Vegetação nativa & 25,0\% & 69,7\% & 22,7\% & 23,7\% \\
\bottomrule
\end{tabular}

\end{table}

\begin{table}[htbp]
\centering
\caption{Correlação de Spearman da aptidão média municipal com o \emph{proxy} de NPP
anual do MODIS, sobre terra manejada, na \textbf{metade de validação} da partição espacial
(Seção~\ref{sec:holdout}). A média municipal usa apenas os pixels dessa metade; entram, assim, os 159
dos 247 municípios que nela têm terra manejada. Uma média parcial é mais ruidosa, o que atenua
$\rho$ e torna a estimativa conservadora.}
\label{tab:spearman-npp}
% GENERATED by tools/gen_thesis_tables.py -- do not edit by hand.
% Edit the source CSV and regenerate; see CLAUDE.md §9.
\begin{tabular}{@{}lrl@{}}
\toprule
Segmento & $\rho$ de Spearman & Significância \\
\midrule
Outras culturas anuais & $+$0,54 & $p < 0{,}001$ \\
Pecuária & $+$0,51 & $p < 0{,}001$ \\
Conservação & $+$0,42 & $p < 0{,}001$ \\
Geração fotovoltaica & $+$0,38 & $p < 0{,}001$ \\
Cana-de-açúcar & $+$0,29 & $p < 0{,}001$ \\
Soja & $+$0,22 & $p = 0{,}004$ \\
Piscicultura & $-$0,30 & $p < 0{,}001$ \\
\bottomrule
\end{tabular}

\end{table}

\chapter{Figuras de diagnóstico}\label{anx:figures}

Este anexo reúne as figuras de diagnóstico citadas nos Capítulos~\ref{ch:results}
e~\ref{ch:discussion}.

\begin{figure}[!hbp]
\centering
\includegraphics[width=0.85\textwidth,height=0.58\textheight,keepaspectratio]{fig_4_2.png}
\caption{Classes de aptidão da FAO (N, S3, S2 e S1) para segmentos representativos
(soja e conservação).}
\label{fig:fao-classes}
\end{figure}

\begin{figure}[htbp]
\centering
\includegraphics[width=0.7\textwidth]{fig_4_6.png}
\caption{Realocação da cultura mais apta sob SSP5-8.5, 2051--2070: células em que muda a cultura de
maior aptidão entre soja, cana-de-açúcar e outras culturas anuais.}
\label{fig:best-crop}
\end{figure}

\begin{figure}[htbp]
\centering
\includegraphics[width=0.7\textwidth]{fig_4_7.png}
\caption{Uso atual da terra de GO~e~DF agregado a 250~m, colorido por função de segmento. É o
retrato espacial da composição tabulada na Seção~\ref{sec:realized-composition}.}
\label{fig:realized-use}
\end{figure}

\begin{figure}[htbp]
\centering
\includegraphics[width=0.7\textwidth]{fig_4_8.png}
\caption{Subutilização da soja (potencial \emph{versus} uso atual): células aptas (S2+), mas
atualmente não ocupadas por soja.}
\label{fig:underused-soy}
\end{figure}

\begin{figure}[htbp]
\centering
\includegraphics[width=\textwidth]{fig_4_13.png}
\caption{Aptidão média municipal para os sete segmentos. Equivalente municipal da
Figura~\ref{fig:atlas-present}, agregado por redução zonal; complementa a Figura~\ref{fig:municipal},
que detalha apenas soja e outras culturas anuais.}
\label{fig:all-segment-suit}
\end{figure}

\section{Seleção do número de zonas ($k$)}\label{sec:k-selection-annex}

\begin{figure}[htbp]
\centering
\includegraphics[width=0.9\textwidth]{fig_4_10.png}
\caption{Índices de validade de agrupamento (silhueta, Davies--Bouldin e estatística de
\emph{gap}) e estabilidade em subamostras (índice de Rand ajustado) para $k$ de 2 a 20, no espaço de
componentes principais descorrelacionado.}
\label{fig:cluster-validity}
\end{figure}

Os índices medidos na metade de validação da partição espacial são os que sustentam a escolha de $k$
(Seção~\ref{sec:zoning}); os da metade de calibração constam ao lado para comparação. Na metade de
validação, a silhueta é quase plana ao longo de todo o intervalo (0,144--0,235), com máximo em
$k=6$; o Davies--Bouldin tem mínimo em $k=6$ (1,400) e varia pouco a partir de $k=7$; e o \emph{gap}
cresce quase monotonicamente e só satisfaz em $k=14$ a regra de um erro-padrão de
\citet[p.~415]{tibshirani_estimating_2001}, segundo a qual se escolhe o menor $k$ tal que
$\mathrm{Gap}(k) \ge \mathrm{Gap}(k+1) - s_{k+1}$. Nenhum desses três índices elege $k=4$.

O que elege $k=4$ (linha tracejada) é o \textbf{ARI de validação cruzada espacial} definido na
Seção~\ref{sec:zoning}: nele, $k=4$ é o máximo claro de todo o intervalo, com 0,793, contra
0,728 em $k=2$, 0,719 em $k=6$ e 0,677 em $k=5$. O ARI de subamostragem concorda: $k=4$ marca
0,986, com o menor desvio-padrão entre réplicas de todo o intervalo ($\pm$0,006; a título de
comparação, $\pm$0,201 em $k=5$ e $\pm$0,176 em $k=6$). Esse índice, porém, mede reprodutibilidade
sob reamostragem, e não validade fora da amostra. Por isso, serve como corroboração, e não como
fundamento.

O painel inteiro varia pouco, e é por isso que a escolha se apoia na reprodutibilidade da
partição em território não usado no ajuste, e não em um ótimo de geometria. O painel cobre $k$ até
20, o que confirma que a escolha não é um artefato de uma faixa truncada.

\section{Decomposição por fator}\label{sec:factor-variance-annex}

\begin{sidewaysfigure}[p]
\centering
\includegraphics[width=0.95\textheight]{fig_4_11.png}
\caption{Parcela de cada fator na variância espacial da aptidão de cada segmento, por
decomposição exata de covariância entre cada fator ponderado e o agregado (as parcelas somam
exatamente~1 por construção). O eixo horizontal está em escala logarítmica simétrica; a cor indica o
grupo temático do fator (clima, laranja; terreno/solo, verde; outros, azul); um contorno preto grosso
marca o fator dominante de cada segmento; a hachura indica contribuição negativa.}
\label{fig:factor-variance}
\end{sidewaysfigure}

A decomposição é feita sobre log-pertinências para os seis segmentos de média geométrica e sobre
pertinências brutas para a conservação, de média aritmética. O eixo horizontal usa uma escala
logarítmica simétrica (limiar linear de $0{,}005$) para manter visíveis as parcelas próximas de zero
sem exagerar sua magnitude real. As contribuições negativas, indicadas por hachura, são um efeito
residual esperado de uma decomposição linear de covariância entre fatores correlacionados, e não um erro.

Terreno e solo dominam as culturas, mas o fator vencedor não é o mesmo em nenhuma delas. Na soja,
a declividade explica 0,93 da variância (argila, 0,07); nas outras culturas anuais, domina o carbono
orgânico do solo (0,66, contra 0,32 da declividade); e na cana-de-açúcar, a água disponível
no solo (0,47), seguida da argila (0,35) e da declividade (0,18). A piscicultura se concentra de forma
quase total na distância à água sazonal (0,91). Nesses quatro segmentos, portanto, a discriminação
espacial é de solo e terreno, mas se reparte de maneira própria em cada um, e não há um fator
único que explique o conjunto.

Os três segmentos restantes não seguem esse padrão. A pecuária é governada pela umidade do solo
(0,97), e a geração fotovoltaica, pela radiação solar (0,71, com o índice de claridade em 0,27). Ambos
são fatores de origem climática, o que distingue esses dois segmentos dos demais. No caso da pecuária,
a umidade do solo descreve o armazenamento de água no perfil, e não um gradiente climático. Se essa
variável for reclassificada como terreno/solo, a parcela desse grupo na pecuária passa de 0,04 para
1,00, e a média climática dos sete segmentos cai de 0,24 para 0,10. Na geração fotovoltaica, por outro
lado, a irradiância \emph{é} um sinal climático espacial real, e a concentração também. A conservação,
por fim, é o segmento com mais fatores do conjunto (nove) e, ainda assim, a conectividade a áreas
protegidas explica sozinha 0,68 da variância, contra 0,15 da rugosidade do terreno, 0,10 do serviço
de carbono e 0,08 da fração de vegetação nativa. É também o único dos sete cujo fator dominante não é
uma variável de superfície, e sim uma camada de governança.

Essa concentração não decorre de viés de ponderação AHP. Nos quatro segmentos dominados por solo e
terreno, ela reflete a saturação geográfica real das pertinências climáticas, confirmada de forma
independente pelo mapa de magnitude espacial da Figura~\ref{fig:scale-mismatch} (argumento
completo na Seção~\ref{sec:atlas-present}). Nos três restantes, reflete a estrutura do próprio
conjunto de fatores de cada segmento.

\section{Descompasso de escala entre clima e terreno}\label{sec:scale-mismatch-annex}

\begin{figure}[htbp]
\centering
\includegraphics[width=0.95\textwidth]{fig_4_12.png}
\caption{Variabilidade espacial local (desvio-padrão em janela móvel, bandas padronizadas)
do clima \emph{versus} terreno/solo, em um raio de 4.750~m (aproximadamente o tamanho da
célula nativa do TerraClimate, $\tfrac{1}{24}^\circ \approx 4{,}6$~km).}
\label{fig:scale-mismatch}
\end{figure}

\begin{table}[htbp]
\centering
\caption{Desvio-padrão local médio (bandas padronizadas) do clima \emph{versus}
terreno/solo, por raio de janela móvel, e razão (terreno/solo)/clima correspondente.}
\label{tab:scale-mismatch}
\begin{tabular}{@{}rrrrr@{}}
\toprule
Raio (m) & $\sigma$ clima & $\sigma$ terreno/solo & Razão & $\log_2$(razão) \\
\midrule
250 & 0,001 & 0,434 & 318 & 8,31 \\
750 & 0,003 & 0,652 & 197 & 7,62 \\
4.750 & 0,022 & 0,899 & 41 & 5,37 \\
\bottomrule
\end{tabular}
\end{table}

O painel do clima (Figura~\ref{fig:scale-mismatch}) é quase uniformemente plano, com variação local
desprezível em toda a extensão do território. O painel de terreno/solo mostra estrutura espacial real,
mais acentuada no norte e no nordeste (transição de relevo do planalto). A razão $\log_2$ entre os dois
painéis é quase uniformemente positiva, o que confirma que terreno e solo dominam a estrutura espacial
fina mesmo na própria escala nativa do clima. O mapa é construído a partir da mesma amostra, limitada
por orçamento computacional (4.475 pontos por raio), que alimenta a Tabela~\ref{tab:scale-mismatch}, e
não de uma composição matricial completa do território. Nessa tabela, o raio de 4.750~m corresponde à
célula nativa do TerraClimate ($\approx$4,6~km); 250~m é a célula nativa do conjunto de atributos; e
750~m reutiliza o raio já usado para a rugosidade do terreno na conservação (Anexo~\ref{anx:derivations}).

\section{Dimensionamento do bloco de validação}\label{anx:holdout}

Esta seção registra as medidas que dimensionaram a partição espacial de calibração e validação
descrita na Seção~\ref{sec:holdout}.

O semivariograma empírico (Tabela~\ref{tab:variograma}) foi estimado para cada cultura, em duas
janelas de distância, sob três definições do campo: a superfície de aptidão ajustada, a presença
observada no MapBiomas e o resíduo entre ambas. O bloco é dimensionado pelo resíduo,
porque é nele que reside a dependência espacial capaz de inflar uma métrica fora da amostra
\citep{roberts_crossvalidation_2017,valavi_blockcv_2019}. O que um bloco precisa separar não é a
suavidade do preditor nem a agregação da cultura, mas a parcela não explicada pelo modelo que,
ainda assim, é compartilhada entre células vizinhas. As outras duas definições são apresentadas como
contraste, pois levariam a dimensionamentos distintos: o alcance da superfície ajustada superestima a
exigência para soja e cana-de-açúcar (155 e 135~km, contra 105~km do resíduo) e a subestima para as
outras culturas anuais (155 contra 175~km).

A pepita é apresentada ao lado de cada alcance porque delimita a sua interpretação. Quando a pepita
se aproxima do patamar, a variância é quase toda ruído na escala de suporte de 250~m, e nenhum
estimador recupera um alcance estrutural. É o caso da presença de outras culturas anuais, cuja
área cultivada cobre menos de 1\% do território: a razão de 0,87 indica que essa linha não
sustenta inferência, ao contrário do resíduo da mesma cultura, com razão de 0,75.

\begin{table}[htbp]
\centering
\caption{Semivariograma empírico das superfícies de cultura sob três enquadramentos
(superfície ajustada, presença observada e resíduo) e duas janelas de distância. ``Compart.''
é a fração da variância espacialmente estruturada ainda compartilhada à distância do bloco
adotado; ``---'' indica alcance não estimável. Subamostra de 6.000 pontos, semente 42.}
\label{tab:variograma}
\small
% GENERATED by tools/gen_thesis_tables.py -- do not edit by hand.
% Edit the source CSV and regenerate; see CLAUDE.md §9.
\begin{tabular}{@{}llrrrr@{}}
\toprule
Cultura & Enquadramento & \makecell[r]{Janela\\(km)} & \makecell[r]{Alcance\\(km)} & \makecell[r]{Pepita/\\patamar} & \makecell[r]{Compart.\\em 130~km} \\
\midrule
Outras c. & presença & 400 & 285 & 0,87 & 100\% \\
 & presença & 800 & --- & 0,73 & 100\% \\
 & resíduo & 400 & 175 & 0,75 & 30\% \\
 & resíduo & 800 & 550 & 0,68 & 66\% \\
 & superfície ajustada & 400 & 155 & 0,75 & 27\% \\
 & superfície ajustada & 800 & 530 & 0,68 & 64\% \\
Soja & presença & 400 & 115 & 0,67 & 7\% \\
 & presença & 800 & 490 & 0,67 & 29\% \\
 & resíduo & 400 & 105 & 0,79 & 11\% \\
 & resíduo & 800 & 150 & 0,79 & 30\% \\
 & superfície ajustada & 400 & 155 & 0,77 & 24\% \\
 & superfície ajustada & 800 & 490 & 0,73 & 55\% \\
Cana & presença & 400 & 165 & 0,70 & 25\% \\
 & presença & 800 & 590 & 0,60 & 63\% \\
 & resíduo & 400 & 105 & 0,78 & 15\% \\
 & resíduo & 800 & 530 & 0,76 & 47\% \\
 & superfície ajustada & 400 & 135 & 0,73 & 20\% \\
 & superfície ajustada & 800 & 530 & 0,67 & 56\% \\
\bottomrule
\end{tabular}

\end{table}

\begin{figure}[htbp]
\centering
\includegraphics[width=\textwidth]{fig_4_15.png}
\caption{Semivariogramas normalizados pelo respectivo patamar, por cultura, na janela de 400~km.
A curva grossa é o resíduo, que dimensiona o bloco; a linha tracejada preta marca o bloco adotado
(130~km), e a pontilhada, o alcance estimado do resíduo. A ordenada inicia próxima de 0,75 porque a
pepita responde por cerca de três quartos do patamar: a maior parte da variância é ruído na escala
de suporte de 250~m, e só o quarto restante é estrutura espacial que um bloco pode separar.}
\label{fig:variograma}
\end{figure}

Os alcances citados acima são da janela de 400~km. Na janela de 800~km, o alcance estimado do
resíduo sobe de 105 para 150~km na soja, de 105 para 530~km na cana-de-açúcar e de 175 para 550~km nas
outras culturas anuais, e a fração compartilhada à distância de 130~km sobe para 30\%, 47\% e
66\%, respectivamente (Tabela~\ref{tab:variograma}). Um alcance que cresce com a janela de
observação indica que não há alcance finito, e sim uma tendência regional de longo alcance. O
dimensionamento foi feito na janela de 400~km porque, a 800~km, as maiores distâncias de defasagem se
aproximam do diâmetro de GO~e~DF. O semivariograma, sustentado então por poucos pares muito distantes,
passa a medir essa tendência, e não a estrutura local. Trata-se, ainda assim, de uma escolha: nenhum
bloco contido no território remove uma tendência que o atravessa por inteiro, e o que a partição de
130~km remove é a estrutura de alcance intermediário.

A Tabela~\ref{tab:split-balance} apresenta a composição por função de uso atual de cada metade, em
uma amostra de 20.000 pixels. A soja, validador primário, fica praticamente equilibrada (razão de 1,01),
e pastagem, vegetação nativa e outras culturas anuais ficam abaixo de 1,4. A cana-de-açúcar é a
exceção, com razão de 4,07: sua área cultivada se concentra no sul de Goiás e, com blocos de
130~km, cai majoritariamente na metade de calibração. Por isso, a linha da cana-de-açúcar nas métricas
de presença deve ser lida à luz dessa limitação. Blocos menores equilibrariam melhor a composição, mas
reduziriam menos a correlação espacial; o compromisso foi resolvido em favor da redução da correlação,
que é a finalidade da partição.

\begin{table}[htbp]
\centering
\caption{Composição por função de uso atual (MapBiomas) de cada metade da partição espacial,
em número de pixels de uma amostra de 20.000 pontos (semente 7, 250~m).}
\label{tab:split-balance}
% GENERATED by tools/gen_thesis_tables.py -- do not edit by hand.
% Edit the source CSV and regenerate; see CLAUDE.md §9.
\begin{tabular}{@{}lrrr@{}}
\toprule
Função de uso atual & Calibração & Validação & Razão \\
\midrule
Outros/sem classe & 251 & 267 & 1,06 \\
Soja & 1.459 & 1.441 & 1,01 \\
Cana-de-açúcar & 362 & 89 & 4,07 \\
Outras culturas anuais & 49 & 57 & 1,16 \\
Pastagem & 4.812 & 4.375 & 1,10 \\
Vegetação nativa & 2.886 & 3.952 & 1,37 \\
\midrule
\textbf{Total} & 9.819 & 10.181 &  \\
\bottomrule
\end{tabular}

\end{table}
